% Job posting: https://ca.linkedin.com/jobs/view/research-engineer-sensor-signal-processing-at-waabi-4466325762
\documentclass[11pt]{article}
\usepackage[left=1in,top=1in,right=1in,bottom=1in]{geometry}
\usepackage[hidelinks]{hyperref}
\usepackage{setspace}
\usepackage[T1]{fontenc}
\usepackage{newtxtext,newtxmath}
\ifdefined\pdfgentounicode
  \input{glyphtounicode}
  \pdfgentounicode=1
\fi
\setstretch{1.1}
\pagenumbering{gobble}
\begin{document}
\pagestyle{empty}

\begin{flushleft}
Nishal Stanislaus Silva, PhD \\
Toronto, ON \\
nishal.silva@hotmail.com \,|\, +1 613 200 2030 \\
\href{https://nishal.xyz}{nishal.xyz} \,|\, \href{https://www.linkedin.com/in/nishal-silva}{linkedin.com/in/nishal-silva}
\end{flushleft}

\vspace{1em}

Dear Hiring Manager,

\vspace{0.5em}

My interest in signal processing started with a time-frequency problem: separating a real onset from noise in a live audio stream. That work became my MSc thesis and an IEEE Asilomar paper (2018) applying the Stockwell (S-)transform to musical onset detection -- a direct exercise in extracting reliable events from a noisy, continuous sensor signal. I have spent the years since turning that kind of signal-processing depth into systems that actually run under real-time constraints: as a postdoctoral researcher at the University of Trento, I owned the full pipeline for a streaming sensor + audio pattern-detection system, from DSP feature engineering through to edge deployment, delivering 14~ms inference on a Raspberry Pi~4 at F1 0.76 -- 74 times faster than the 1038~ms baseline it replaced (J. Audio Eng. Soc., 2026). That is the kind of role I am looking for next: one where signal-processing rigor has to survive contact with a real-time deployment target, which is how I read the Research Engineer, Sensor Signal Processing position at Waabi.

\vspace{0.5em}

Beyond single-stream processing, I have direct experience fusing multiple sensor modalities in real time. At McGill University's IDMIL/CIRMMT, I built a self-powered smart electric guitar instrumented with an IMU and audio interface, then fused the two streams through a hierarchical finite-state machine and an LSTM gesture model, reaching F1 0.78 on a 500-event corpus (IEEE IS2, 2025). That project required the same discipline sensor signal processing for autonomous systems demands: synchronizing heterogeneous sensor streams, profiling compute and power budgets on embedded hardware, and making the fusion logic robust to real-world noise. My DSP toolkit (STFT, MFCC/LFCC/GFCC, CQT, onset and pitch tracking) and my C++17 systems background, sharpened over 5.5 years building embedded computer-vision and IoT pipelines at MAS Holdings, round out a profile built for shipping signal-processing code that has to run, not just publish.

\vspace{0.5em}

What I want to be upfront about: I have not worked with radar or lidar, and I have no autonomous-vehicle-specific sensor experience. What transfers directly is everything underneath the sensor modality -- the signal-processing fundamentals, the multi-sensor fusion architecture, and the embedded real-time engineering discipline. The specific physics and data characteristics of automotive-grade sensors are the part I would need to learn quickly, and I am confident that learning curve is short given the adjacent depth above.

\vspace{0.5em}

I am a Canadian permanent resident, based in Toronto, and available to work without sponsorship. I would welcome the chance to discuss how this background could contribute to Waabi's sensor signal processing team.

\vspace{1em}

Sincerely, \\
Nishal Stanislaus Silva

\end{document}
